\section{Бегущий огонек}
\begin{verbatim}
					Начальное состояние: выкл
					Движение: От крайних к середине
					Порядок: 09-18-27-36-45-36-27-18-09
\end{verbatim}
Код прошивки:
\inputminted{cpp}{z1/main.cpp}
i - это число от 0 до 9, которое используется для прямого и обратного "хода"\ светодиода левой половины и это отзеркаливается на правую половину.
\section{Пульсар}
\begin{verbatim}
					Начальное состояние: выкл
					Транзистор: Полевой
					Смена состояний: Увеличение до середны/сброс/Увеличение max
\end{verbatim}
Код прошивки:
\inputminted{cpp}{z2/main.cpp}
to\_half является флагом для прохода до 127, потом сбрасывается и счётчик идёт до 255.
\section{Ночной светильник}
\begin{verbatim}
					Начальное состояние: вкл
					Входной сигнал: фоторезистор
					Делитель напряжения: 10 кОм
					Смена состояний: Уменьшение в зависимости
						от значения на датчике
\end{verbatim}
Код прошивки:
\inputminted{cpp}{z3/main.cpp}
Получаемый с пина A0 сигнал от фоторезистора в аудитори в среднем имеет значение 890-910, при полном закрытии он выдаёт 160-350, а при свете фонарика он максымум выдаст 1000-1010. Для максимальной и минимальной яркости светодиода используется вычетание и проверки значения с датчика, а также инвертируется для того чтобы яркость светодиода уменьшалась по мере увеличения значения с датчика.
\section{Азбука Морзе}
\begin{verbatim}
					Кнопка: Срабатывание по нажатию
\end{verbatim}
Файл melody.h:
\inputminted{cpp}{z4/melody.h}
Код прошивки:
\inputminted{cpp}{z4/main.cpp}
Первая кнопка переключает мелодию, также останавливает текущую мелодию и сбрасывает указатель. В "melody.h"\ хранятся мелодии в виде нот, то есть, в массиве где записана частота ноты и её задержка. Флаги "first\_selected"\ и "running"\ отвечают за выбор и воспроизведение мелодии соответсвенно. В зависимости от выбранной мелодии загорается первый или второй светодиод, а также третий светодиод мигает когда звучит не нулевая частота.
\section{Терменвокс}
\begin{verbatim}
					При падении освещенности звук: Уменьшается
					Звуковой сигнал: 10 раз в сек.
					Частоты нот: Малая октава
\end{verbatim}
Код прошивки:
\inputminted{cpp}{z5/main.cpp}
"frequency"\ хранит частоту, которую мы получаем с термистора путём преобразования его диапазона в диапазон частот Малой октавы. На термисторе выбран диапазон 400 - 500 т.к. в аудитории он остывает и стремится к 500, а при нагреве пальцами стремится к 400.
\section{Мерзкое пианино}
\begin{verbatim}
					Подключение кнопок: Без резистора
					Входной сигнал: массив
					Частоты нот: Малая октава
\end{verbatim}
Код прошивки:
\inputminted{cpp}{z6/main.cpp}
"index"\ - как указатель в цикле, проходит по значениям от 0 до 4. в "b"\ хранятся предыдущие состояния кнопок нажата/не нажата, а в "m"\ хранятся частоты. в "btn"\ записывается текщее состояние текущей кнопки и сравнивается с предыдущим её состоянием. Для реализации без резистора была добавлена настройка пинов на встроенный резистор, чтобы небыло помех.